% GENERATED FILE. Edit canonical lessons, then run npm run generate:books.
% canonical-source-hash: fnv1a64:5f460afd51caa639
% canonical-lessons: AR-C520-itada, AR-C520-munawaba, AR-C520-ijaza, AR-C520-mahamma, AR-C520-naffadha

\chapter{Hours and tasks}
\label{ch:ar-hours-and-tasks}
\hlchaptermodality{520}

\begin{chapteropening}
\textbf{By the end of this chapter:} \emph{I can talk about a shift, time off and a task, and say that someone carried out a plan or got used to the work.}

Complete the last lesson of chapter 520: I can talk about a shift, time off and a task, and say that someone carried out a plan or got used to the work.
\end{chapteropening}

\section[\texorpdfstring{\ar{اعتاد} (iʿtāda)}{اعتاد (iʿtāda)}]{\ar{اعتاد} (iʿtāda) — to get used to}

\label{lesson:AR-C520-itada}



\begin{quote}
Before the new one: say the Arabic for a trainee, then the Arabic for to practise.
\end{quote}

\subsection*{You'll want to know: \ar{اعتاد}}
\textbf{\ar{اعتاد}} — \emph{iʿtāda} — \textquotedblleft{}to get used to\textquotedblright{}.

\begin{quote}
\ar{اعتادت} \ar{الممرضة} \ar{على} \ar{العمل} \ar{في} \ar{الليل}.
\end{quote}

\begin{quote}
\emph{iʿtādat al-mumarriḍa ʿalā al-ʿamal fī al-layl.}
\end{quote}

\begin{quote}
The nurse got used to working at night.
\end{quote}

Present: \ar{يعتاد}, \emph{yaʿtād}.

It takes \ar{على}, or a direct object.

\begin{cousinweb}
\textbf{\ar{اعتاد}} (\emph{iʿtāda}) is built on the native Arabic root \emph{ʿ-w-d}, \textquotedblleft{}to return\textquotedblright{} (Lane). ʿĀda, a habit, which you know, is from the same root: a thing that keeps coming back.
\end{cousinweb}

\begin{grammarlens}[title={What you've built}]
One more word for describing what you lived and felt.
\end{grammarlens}

\subsection*{Your turn}
\begin{itemize}
\raggedright
  \item \emph{Say it:} \emph{iʿtāda}
  \item \emph{Say it:} \emph{iʿtāda}, once more
  \item \emph{Say it:} \emph{iʿtādat al-mumarriḍa ʿalā al-ʿamal fī al-layl.} — The nurse got used to working at night.
  \item \emph{Recall:} say the Arabic for a trainee, then the Arabic for to practise, then say \emph{iʿtāda} again
\end{itemize}

\begin{tcolorbox}[breakable,colback=teal!4,colframe=teal!35!black,title={Before you move on}]
What does \ar{اعتاد} mean? (\textquotedblleft{}To get used to\textquotedblright{}.) Say it once more.
\end{tcolorbox}
\section[\texorpdfstring{\ar{مناوبة} (munāwaba)}{مناوبة (munāwaba)}]{\ar{مناوبة} (munāwaba) — a shift, a turn on duty}

\label{lesson:AR-C520-munawaba}



\begin{quote}
Before the new one: say the Arabic for to practise, then the Arabic for to get used to.
\end{quote}

\subsection*{You'll want to know: \ar{مناوبة}}
\textbf{\ar{مناوبة}} — \emph{munāwaba} — \textquotedblleft{}a shift, a turn on duty\textquotedblright{}.

\begin{quote}
\ar{عند} \ar{الممرضة} \ar{مناوبة} \ar{في} \ar{الليل}.
\end{quote}

\begin{quote}
\emph{ʿinda al-mumarriḍa munāwaba fī al-layl.}
\end{quote}

\begin{quote}
The nurse has a night shift.
\end{quote}

\begin{cousinweb}
\textbf{\ar{مناوبة}} (\emph{munāwaba}) is built on the native Arabic root \emph{n-w-b}, \textquotedblleft{}to come in turn\textquotedblright{} (Lane). Munāwaba is the verbal noun of Form III: taking turns.
\end{cousinweb}

\begin{grammarlens}[title={What you've built}]
One more word for describing what you lived and felt.
\end{grammarlens}

\subsection*{Your turn}
\begin{itemize}
\raggedright
  \item \emph{Say it:} \emph{munāwaba}
  \item \emph{Say it:} \emph{munāwaba}, once more
  \item \emph{Say it:} \emph{ʿinda al-mumarriḍa munāwaba fī al-layl.} — The nurse has a night shift.
  \item \emph{Recall:} say the Arabic for to practise, then the Arabic for to get used to, then say \emph{munāwaba} again
\end{itemize}

\begin{tcolorbox}[breakable,colback=teal!4,colframe=teal!35!black,title={Before you move on}]
What does \ar{مناوبة} mean? (\textquotedblleft{}A shift, a turn on duty\textquotedblright{}.) Say it once more.
\end{tcolorbox}
\section[\texorpdfstring{\ar{إجازة} (ijāza)}{إجازة (ijāza)}]{\ar{إجازة} (ijāza) — time off, leave}

\label{lesson:AR-C520-ijaza}



\begin{quote}
Before the new one: say the Arabic for to get used to, then the Arabic for a shift.
\end{quote}

\subsection*{You'll want to know: \ar{إجازة}}
\textbf{\ar{إجازة}} — \emph{ijāza} — \textquotedblleft{}time off, leave\textquotedblright{}.

\begin{quote}
\ar{أخذت} \ar{أمي} \ar{إجازة} \ar{من} \ar{العمل}.
\end{quote}

\begin{quote}
\emph{akhadhat ummī ijāza min al-ʿamal.}
\end{quote}

\begin{quote}
My mother took time off work.
\end{quote}

It also means a permit or a licence.

\begin{cousinweb}
\textbf{\ar{إجازة}} (\emph{ijāza}) is built on the native Arabic root \emph{j-w-z}, \textquotedblleft{}to pass, to go through\textquotedblright{} (Lane). Jawāz, a passport, which you know, is from the same root: both let you pass.
\end{cousinweb}

\begin{grammarlens}[title={What you've built}]
One more word for describing what you lived and felt.
\end{grammarlens}

\subsection*{Your turn}
\begin{itemize}
\raggedright
  \item \emph{Say it:} \emph{ijāza}
  \item \emph{Say it:} \emph{ijāza}, once more
  \item \emph{Say it:} \emph{akhadhat ummī ijāza min al-ʿamal.} — My mother took time off work.
  \item \emph{Recall:} say the Arabic for to get used to, then the Arabic for a shift, then say \emph{ijāza} again
\end{itemize}

\begin{tcolorbox}[breakable,colback=teal!4,colframe=teal!35!black,title={Before you move on}]
What does \ar{إجازة} mean? (\textquotedblleft{}Time off, leave\textquotedblright{}.) Say it once more.
\end{tcolorbox}
\section[\texorpdfstring{\ar{مهمة} (mahamma)}{مهمة (mahamma)}]{\ar{مهمة} (mahamma) — a task, a mission}

\label{lesson:AR-C520-mahamma}



\begin{quote}
Before the new one: say the Arabic for a shift, then the Arabic for time off.
\end{quote}

\subsection*{You'll want to know: \ar{مهمة}}
\textbf{\ar{مهمة}} — \emph{mahamma} — \textquotedblleft{}a task, a mission\textquotedblright{}.

\begin{quote}
\ar{عندي} \ar{مهمة} \ar{صعبة} \ar{اليوم}.
\end{quote}

\begin{quote}
\emph{ʿindī mahamma ṣaʿba al-yawm.}
\end{quote}

\begin{quote}
I have a difficult task today.
\end{quote}

Plural: \ar{مهام}, \emph{mahāmm}.

\begin{cousinweb}
\textbf{\ar{مهمة}} (\emph{mahamma}) is built on the native Arabic root \emph{h-m-m}, \textquotedblleft{}to care about, to be concerned\textquotedblright{} (Lane). Muhimm, important, which you know, is from the same root.
\end{cousinweb}

\begin{grammarlens}[title={What you've built}]
One more word for describing what you lived and felt.
\end{grammarlens}

\subsection*{Your turn}
\begin{itemize}
\raggedright
  \item \emph{Say it:} \emph{mahamma}
  \item \emph{Say it:} \emph{mahamma}, once more
  \item \emph{Say it:} \emph{ʿindī mahamma ṣaʿba al-yawm.} — I have a difficult task today.
  \item \emph{Recall:} say the Arabic for a shift, then the Arabic for time off, then say \emph{mahamma} again
\end{itemize}

\begin{tcolorbox}[breakable,colback=teal!4,colframe=teal!35!black,title={Before you move on}]
What does \ar{مهمة} mean? (\textquotedblleft{}A task, a mission\textquotedblright{}.) Say it once more.
\end{tcolorbox}
\section[\texorpdfstring{\ar{نفذ} (naffadha)}{نفذ (naffadha)}]{\ar{نفذ} (naffadha) — to carry out}

\label{lesson:AR-C520-naffadha}



\begin{quote}
Before the new one: say the Arabic for time off, then the Arabic for a task.
\end{quote}

\subsection*{You'll want to know: \ar{نفذ}}
\textbf{\ar{نفذ}} — \emph{naffadha} — \textquotedblleft{}to carry out\textquotedblright{}.

\begin{quote}
\ar{نفذ} \ar{العمال} \ar{الخطة} \ar{بسرعة}.
\end{quote}

\begin{quote}
\emph{naffadha al-ʿummāl al-khuṭṭa bi-surʿa.}
\end{quote}

\begin{quote}
The workers carried out the plan quickly.
\end{quote}

Present: \ar{ينفذ}, \emph{yunaffidh}.

\begin{cousinweb}
\textbf{\ar{نفذ}} (\emph{naffadha}) is built on the native Arabic root \emph{n-f-dh}, \textquotedblleft{}to pass through, to go through\textquotedblright{} (Lane). Form II naffadha is to make a plan go through.
\end{cousinweb}

\begin{grammarlens}[title={What you've built}]
One more word for describing what you lived and felt.
\end{grammarlens}

\subsection*{Your turn}
\begin{itemize}
\raggedright
  \item \emph{Say it:} \emph{naffadha}
  \item \emph{Say it:} \emph{naffadha}, once more
  \item \emph{Say it:} \emph{naffadha al-ʿummāl al-khuṭṭa bi-surʿa.} — The workers carried out the plan quickly.
  \item \emph{Recall:} say the Arabic for time off, then the Arabic for a task, then say \emph{naffadha} again
\end{itemize}

\begin{tcolorbox}[breakable,colback=teal!4,colframe=teal!35!black,title={Before you move on}]
What does \ar{نفذ} mean? (\textquotedblleft{}To carry out\textquotedblright{}.) Say it once more.
\end{tcolorbox}
